\subsection{Tensor products, tensor spaces, exterior algebra}

We keep the notation of 7.6. Moreover, we denote by L a commutative field endowed with a structure of K-algebra of finite dimension, and we call vector bundle a vector bundle over L, with base B and locally finite rank.

7.9.1. Suppose that \(\mathbf{I}_{-} = \emptyset\) . If \(\mathcal{V}\) and \(\mathcal{V}'\) are two families indexed by I of finite-dimensional vector spaces over L, we denote by \(\tau(\mathcal{V})\) the tensor product of the \(V_{i}\) for \(i \in \mathbf{I}\) (A, II, p. 71) and if \(\mathbf{f} \in \operatorname{Hom}(\mathcal{V}, \mathcal{V}')\) , we set \(\tau(\mathbf{f}) = \otimes f_{i}\) . One thus defines a vector functor over L in finite

dimension, and if \(\mathcal{M} = (\mathrm{M}_{i})_{i\in I}\) is a family of vector bundles, we denote the vector bundle \(\tau(\mathcal{M})\) by \(\bigotimes_{i\in I}\mathrm{M}_{i}\) and call it the tensor product (over L) of the \(\mathrm{M}_{i}\) .

If \(s_i\) is a section of \(\mathbf{M}_i\) over the open set U of B (for \(i \in I\) ), the map \(b \mapsto \bigotimes_{i \in I} s_i(b)\) is a section of \(\bigotimes_{i \in I} \mathbf{M}_i\) , denoted \(\bigotimes_{i \in I} s_i\) . The map \((s_i)_{i \in I} \mapsto \bigotimes_{i \in I} s_i\) is multilinear over the ring \(\mathcal{C}^r(\mathrm{U}; \mathrm{L})\) .

7.9.2. The canonical isomorphisms defined in Alg., chap. II provide isomorphisms of vector functors. It follows from 7.6.4 that there are isomorphisms of vector bundles. For example, we have canonical isomorphisms:

\[
\begin{array}{c} (\mathrm{M} _ {1} \oplus \mathrm{M} _ {2}) \otimes \mathrm{M} _ {3} \longrightarrow (\mathrm{M} _ {1} \otimes \mathrm{M} _ {3}) \oplus (\mathrm{M} _ {2} \otimes \mathrm{M} _ {3}) \\ \mathrm{M} _ {1} ^ {*} \otimes \mathrm{M} _ {2} \longrightarrow \mathcal {L} (\mathrm{M} _ {1}; \mathrm{M} _ {2}) \end{array}
\]

etc.

7.9.3. Let M be a vector bundle and let I and J be two disjoint finite sets. The tensor bundle \(\mathsf{T}_{\mathbf{J}}^{\mathbf{I}}(\mathbf{M})\) is defined as the tensor product \(\bigotimes_{\alpha\in\mathbf{I}\cup\mathbf{J}}\mathbf{M}_{\alpha}\) , where \(M_{\alpha}=M\) if \(\alpha\in I\) , and \(M_{\alpha}=M^{*}\) if \(\alpha\in J\) , \(M^{*}\) denotes the dual of M (A, III, p. 63). The fiber \(\mathsf{T}_{\mathbf{J}}^{\mathbf{I}}(\mathbf{M})_{b}\) of this bundle at a point b is equal to the tensor space \(\mathsf{T}_{\mathbf{J}}^{\mathbf{I}}(\mathbf{M}_{b})\) defined in Alg., loc. cit. When \(I=\{1,\ldots,p\}\) and \(J=\{p+1,\ldots,p+q\}\) , we write \(\mathsf{T}_{q}^{p}(\mathbf{M})\) instead of \(\mathsf{T}_{\mathbf{J}}^{\mathbf{I}}(\mathbf{M})\) ; we have

\[
\mathbf {T} _ {q} ^ {p} (\mathbf {M}) = (\bigotimes^ {p} \mathbf {M}) \otimes (\bigotimes^ {q} \mathbf {M} ^ {*}).
\]

The data of a total order on I and on J defines a canonical isomorphism of \(\mathsf{T}_{\mathbf{J}}^{\mathbf{I}}(M)\) onto \(\mathsf{T}_{q}^{p}(M)\) .

7.9.4. If I (resp. J) is the union of two disjoint subsets I' and I'' (resp. J' and J''), \(\mathsf{T}_{\mathbf{J}}^{\mathbf{I}}(M)\) is canonically identified with the tensor product \(\mathsf{T}_{\mathbf{J}'}^{\mathbf{I}'}(M) \otimes \mathsf{T}_{\mathbf{J}''}^{\mathbf{I}''}(M)\) . In particular, if s' (resp. s'') is a section of \(\mathsf{T}_{\mathbf{J}'}^{\mathbf{I}'}(M)\) (resp. of \(\mathsf{T}_{\mathbf{J}''}^{\mathbf{I}''}(M)\) ), the tensor product s' \(\otimes\) s'' is identified with a section of \(\mathsf{T}_{\mathbf{J}}^{\mathbf{I}}(M)\) .

The dual of \(\mathsf{T}_{\mathbf{J}}^{\mathbf{I}}(\mathbf{M})\) is identified with \(\mathsf{T}_{\mathbf{I}}^{\mathbf{J}}(\mathbf{M})\) .

7.9.6. Let \(i \in I\) and \(j \in J\) . For every \(b \in B\) , the contraction homomorphism of indices \(i\) and \(j\) is defined (cf. Alg., loc. cit.); it is a homomorphism \((c_j^i)_b: \mathsf{T}_{\mathbf{J}}^{\mathbf{I}}(M_b) \to \mathsf{T}_{\mathbf{J}-\{j\}}^{\mathbf{I}-\{i\}}(M_b)\) . The collection of the \((c_j^i)_b\) defines a morphism of vector bundles

\[
c _ {j} ^ {i}: \mathsf {T} _ {\mathbf {J}} ^ {\mathbf {I}} (\mathbf {M}) \to \mathsf {T} _ {\mathbf {J} - \{j \}} ^ {\mathbf {I} - \{i \}} (\mathbf {M}),
\]

still called contraction of indices \(i\) and \(j\) . One defines similarly the contraction of the indices \(i_{1},\ldots,i_{k}\) of I with the indices \(j_{1},\ldots,j_{k}\) of J.

7.9.7. Let \(d\) be an integer \(\geqslant 0\) ; let V and V' be two finite-dimensional vector spaces over L and \(f \in \mathrm{Hom}_{\mathrm{L}}(\mathrm{V}, \mathrm{V}')\) ; set \(\lambda_{d}(\mathrm{V}) = \bigwedge^{d}(\mathrm{V})\) and \(\lambda_{d}(f) = \bigwedge^{d}(f)\) . One thus defines a vector functor on L in finite dimension and, if M is a vector bundle, one denotes the vector bundle \(\lambda_{d}(\mathbf{M})\) by \(\bigwedge^{d}(\mathbf{M})\), calling it the d-th exterior power of M (over L) .

The canonical map from \(\otimes^{d}\mathbf{V}\) onto \(\bigwedge^{d}(\mathbf{V})\) defines a morphism of vector functors, whence a morphism, said to be canonical, from \(\otimes^{d}\mathbf{M}\) to \(\bigwedge^{d}(\mathbf{M})\) . This morphism is surjective.

The canonical isomorphisms of the space of alternating d-multilinear maps onto the space \(\bigwedge^{d}(\mathrm{V}^{*})\) or onto \((\bigwedge^{d}(\mathrm{V}))^{*}\) provide isomorphisms, said to be canonical, of the vector bundle \(\operatorname{Alt}^{d}(\mathbf{M};\mathbf{L})\) of alternating d-multilinear maps over L, onto the vector bundle \(\bigwedge^{d}(\mathbf{M}^{*})\) or onto \((\bigwedge^{d}(\mathbf{M}))^{*}\) .

7.9.8. Now set \(\lambda(V) = \bigwedge(V)\) and \(\lambda(f) = \bigwedge(f)\) ; one thus again defines a vector functor on \(L\) in finite dimension. The vector bundle \(\lambda(M)\) is denoted \(\bigwedge(M)\) . Its fiber \(\bigwedge(M)_b\) at \(b \in B\) is the exterior algebra (over \(L\) ) of the fiber \(M_b\) . The vector bundle \(\bigwedge(M)\) is a locally trivial vector bundle of algebras.

The definitions and properties of interior products given in Alg., chap. III, 3rd ed., § 10 extend immediately to sections of vector bundles \(\bigwedge(\mathbf{M})\) and \(\bigwedge(\mathbf{M}^{*})\) (cf. also formulas (1) to (7) of nos. 7.8.2 to 7.8.4).

7.9.9. Let M be a vector bundle. For every integer n, let \(B_{n}\) be the (open) set of points \(b \in B\) such that the dimension (over L) of \(M_{b}\) is equal to n. For \(b \in B_{n}\) , set \(N_{b} = \bigwedge^{n}(M_{b})\) and let N be the disjoint union of the \(N_{b}\) for \(b \in B\) . There exists on N a vector bundle structure, and only one, such that the evident map from \(N|B_{n}\) to \(\bigwedge^{n}(M)|B_{n}\) is an isomorphism for every n. Endowed with this structure, the vector bundle N, of rank one at each point, is denoted det(M).

